\documentclass[11pt]{article}

\usepackage[margin=1in]{geometry}
\usepackage{amsmath,amssymb,amsthm}
\usepackage{array,booktabs}
\usepackage{enumitem}
\usepackage[hidelinks]{hyperref}

\theoremstyle{plain}
\newtheorem{observation}{Observation}
\newtheorem{proposition}{Proposition}
\newtheorem{lemma}{Lemma}
\newtheorem{corollary}{Corollary}
\newtheorem{theorem}{Theorem}
\newtheorem*{question}{Open question}
% The lettered series records the no-go results of the 2026-06-10 parallel
% research run; the letters A--H are stable names used by OPEN.md and the
% probe files in experiments/gold/.
\newtheorem{nogo}{Theorem}
\renewcommand{\thenogo}{\Alph{nogo}}
\theoremstyle{definition}
\newtheorem{definition}{Definition}
\newtheorem{remark}{Remark}

\newcommand{\No}{\mathbf{No}}
\newcommand{\On}{\mathrm{On}_2}
\newcommand{\F}{\mathbb{F}}
\newcommand{\Tr}{\operatorname{Tr}}
\newcommand{\Arf}{\operatorname{Arf}}
\newcommand{\nimmul}{\otimes}
\newcommand{\Cl}{\mathrm{Cl}}
\newcommand{\Sp}{\operatorname{Sp}}
\newcommand{\Orth}{\mathrm{O}}
\newcommand{\AO}{\mathrm{AO}}
\newcommand{\AGL}{\operatorname{AGL}}
\newcommand{\GL}{\operatorname{GL}}
\newcommand{\Aut}{\operatorname{Aut}}
\newcommand{\Inn}{\operatorname{Inn}}
\newcommand{\Out}{\operatorname{Out}}
\newcommand{\rad}{\operatorname{rad}}
\newcommand{\wt}{\operatorname{wt}}
\newcommand{\mex}{\operatorname{mex}}
\newcommand{\Stab}{\operatorname{Stab}}

\title{Draft: game-built quadratic forms in the nimber Clifford backend}
\author{a9lim}
\date{June 2026}

\begin{document}
\maketitle

\begin{abstract}
This is a consolidated experimental note about the part of \texttt{ogdoad} that
seems mathematically worth isolating. Conway games under disjunctive sum do not
form a scalar ring, so the project does not construct ``Clifford algebras over
all games''. Instead it builds Clifford algebras over field-like game-value
worlds, with emphasis on the characteristic-$2$ nimber backend, whose engine
keeps the quadratic data $q_i=e_i^2$ and the polar data
$b_{ij}=e_ie_j+e_je_i$ independent and computes trace-Arf invariants of
quadratic forms over finite nimber subfields.

The experimental bridge is this: the Gold trace forms
\[
  Q_a(x)=\Tr(x\,x^{2^a})
\]
are standard quadratic forms over fields of characteristic $2$, but in the
nimber setting they are also composites of combinatorial-game operations:
nim-product via Turning Corners, Frobenius squaring via the diagonal product,
and XOR via disjunctive sum. Their Arf invariant has the classical zero-count
interpretation, so if a natural game had $P$-positions exactly
$\{x:Q_a(x)=0\}$, Arf would be the sign of the second-player win-bias. The
note does not claim such a game has been found.

Beyond the construction, the note synthesizes three threads of work on the
open game-semantics question. First, the obstruction landscape: normal-play
sums and Winning-Ways coin-turning realize only \emph{linear} $P$-sets
(the lexicode phenomenon is the rich solved case of that linear ceiling), and
a family of no-go theorems eliminates symmetric, commutative, oracle-bounded,
and translation-equivariant architectures above it. Second, an
extraspecial-group reframing: a characteristic-$2$ quadratic form \emph{is} a
central extension $1\to\mathbb{Z}/2\to E\to V\to 1$ whose squaring map is $Q$,
which forces the quadratic datum to come from a noncommutative source and
yields a naturality screen strictly between the frame-blind no-go and
tautological evaluation. Third, the constructive boundary: within every tested
bounded-access class, exact realizers of $\{Q=0\}$ exist but are forced
clocks; one charge-counting construction (\textsc{echo}-ko) is exact on
the verified rank-$4$ bent instances and near-exact at low rank, while
remaining decision-live. A formalized naturality
criterion makes the residual question precise: does a decision-nondegenerate
fixed uniform rule realize a Gold quadric?
\end{abstract}

\paragraph{Claim levels.}
Following the repository convention, statements below are labelled
\emph{standard math} (external facts with citations), \emph{implemented and
tested} (backed by tests, examples, or experiments in this checkout),
\emph{interpretation} (conditional bridges), or \emph{open}. Several results
of the 2026-06-10 parallel research run are machine-verified only at stated
small parameters; each such result carries its verification scope inline.
One construction claim (\textsc{echo}-\textsc{fifo}+dummy,
Section~\ref{sec:constructions}) is explicitly \emph{unverified} and is not
relied on anywhere.

\section{Scope}\label{sec:scope}

A short partizan game is built recursively as $\{L\mid R\}$, and games form a
partially ordered abelian group under disjunctive sum~\cite{Conway76,WW}. They
do not form a ring: Conway multiplication is defined on the number subclass,
not on arbitrary games. Since a Clifford algebra needs a commutative scalar
ring, the phrase ``Clifford over games'' has to be interpreted through scalar
subworlds:

\begin{center}
\begin{tabular}{@{}>{\raggedright\arraybackslash}p{0.18\linewidth}
>{\raggedright\arraybackslash}p{0.28\linewidth}
>{\raggedright\arraybackslash}p{0.45\linewidth}@{}}
\toprule
world & algebraic structure & role here \\
\midrule
$\No$ & real-closed field, char $0$ & ordinary real Clifford theory with surreal scalars \\
$\No[i]$ & algebraically closed field, char $0$ & ordinary complex Clifford theory \\
finite nimber subfields & finite fields of char $2$ & characteristic-$2$ Clifford and Arf experiments \\
\bottomrule
\end{tabular}
\end{center}

The implemented \texttt{Nimber} scalar is finite: \texttt{u128} realizes
$\F_{2^{128}}$, containing the nimber subfields $\F_{2^m}$ for
$m=1,2,4,\ldots,128$. It is not the whole proper-class field $\On$.

\section{The char-2 Clifford datum}\label{sec:datum}

In characteristic $2$, a quadratic form is not determined by its polar form.
For basis vectors the Clifford relations are
\[
  e_i^2=q_i,\qquad e_i e_j+e_j e_i=b_{ij}.
\]
The polar form is alternating, so $b_{ii}=0$, but $q_i$ can be nonzero. A
faithful characteristic-$2$ implementation must store $q$ and $b$ separately.
This is the main implementation discipline behind the nimber backend, and --
as Section~\ref{sec:extraspecial} makes precise -- it is also the
group-theoretic content of the whole note: $q$ and $b$ are the independent
central data (squares and commutators) of an extraspecial-type extension.

For a nonsingular quadratic form over $\F_2$ with symplectic basis
$(a_k,b_k)$, the Arf invariant is
\[
  \Arf(Q)=\sum_k Q(a_k)Q(b_k)\in \F_2.
\]
It classifies nonsingular binary quadratic spaces at fixed dimension. Over a
finite characteristic-$2$ field, the corresponding value lies in
$F/\wp(F)$, with $\wp(y)=y^2+y$; for finite fields this quotient is read by the
field trace to $\F_2$.

\begin{remark}
Ovsienko~\cite{Ovsienko16} identifies a precise bridge between Dickson's
classification of quadratic forms over $\F_2$ and the classical classification
of real Clifford algebras. We use that paper as motivation for the Arf/Clifford
connection, not as a blanket classification theorem for all Clifford algebras
over finite nimber fields. In the code, rank and radical data are reported
alongside Arf, and degenerate forms are not collapsed to a single Arf bit.
\end{remark}

The small finite-field check used throughout is:
\[
  q=(\ast2,\ast2),\ b_{01}=\ast1 \quad\text{has Arf }1,
\]
whereas
\[
  q=(\ast2,\ast3),\ b_{01}=\ast1 \quad\text{has Arf }0.
\]
The polar entry is part of the form.

\section{Gold forms from game operations}\label{sec:gold}

Nim-addition is XOR, and XOR is the disjunctive sum of impartial game values.
Nim-multiplication is also game-theoretic: Conway's Turning-Corners product has
the mex recurrence
\begin{equation}\label{eq:tc}
  x\nimmul y
  =
  \mex
  \left\{
    (i\nimmul y)\oplus(x\nimmul j)\oplus(i\nimmul j):
    0\leq i<x,\ 0\leq j<y
  \right\}.
\end{equation}
The repo implements this recurrence directly for small arguments and checks it
against the fast algebraic nim product.

Thus the following operations are game-realizable at the level of nimber
values:
\[
  x\oplus y,\qquad x\nimmul y,\qquad x\mapsto x^2=x\nimmul x,\qquad
  \Tr(x)=x+x^2+\cdots+x^{2^{m-1}}.
\]

For $F=\F_{2^m}$ and $a\geq 1$, define the Gold trace form
\[
  Q_a(x)=\Tr\!\left(x\,x^{2^a}\right)=\Tr\!\left(x^{1+2^a}\right).
\]
Since Frobenius $x\mapsto x^{2^a}$ is additive, $Q_a$ is a quadratic form over
$\F_2$. In the nimber backend, it is also a composite of the game operations
above. Throughout, ``the instance $(m,a)$'' means $Q_a$ on $\F_{2^m}$, and
``$(m,a,\lambda)$'' the scaled component $Q_{\lambda,a}(x)=\Tr(\lambda
x^{1+2^a})$.

\begin{observation}\label{obs:built}
The Gold forms are not merely defined on a field of game values. In the finite
nimber fields tested here, they are built from nim/game operations: Turning
Corners for multiplication, diagonal Turning Corners for Frobenius squaring,
and disjunctive sum for XOR and trace.
\end{observation}

This observation is modest but useful: it gives a concrete quadratic form at
the interface between impartial-game arithmetic and characteristic-$2$
quadratic-form theory.

\subsection{Rank check}

The polar rank of Gold forms is known. In the tested regime
$m=2^k$ and $1\leq a\leq k$, the general formula
\[
  \operatorname{rank} B_{Q_a}=m-\gcd(2a,m)
\]
specializes to $m-2\gcd(a,m)$ because $m/\gcd(a,m)$ is even. The classifier
reproduces the general formula in all tested cases. Representative rows are
shown in Table~\ref{tab:gold}.

\begin{table}[h]
\centering
\begin{tabular}{rrccccc}
\toprule
field & $a$ & $Q$ & $\Arf$ & type & rank & formula \\
\midrule
$\F_{2^{4}}$  & $1$ & $\Tr(x^3)$    & $1$ & $O^-$ & $2$  & $2$ \\
$\F_{2^{8}}$  & $1$ & $\Tr(x^3)$    & $1$ & $O^-$ & $6$  & $6$ \\
$\F_{2^{8}}$  & $2$ & $\Tr(x^5)$    & $1$ & $O^-$ & $4$  & $4$ \\
$\F_{2^{16}}$ & $1$ & $\Tr(x^3)$    & $1$ & $O^-$ & $14$ & $14$ \\
$\F_{2^{16}}$ & $4$ & $\Tr(x^{17})$ & $1$ & $O^-$ & $8$  & $8$ \\
$\F_{2^{32}}$ & $1$ & $\Tr(x^3)$    & $1$ & $O^-$ & $30$ & $30$ \\
\bottomrule
\end{tabular}
\caption{Representative output of \texttt{experiments/trace\_form\_arf.py}.
The script checks all $15$ cases with $m=2,4,8,16,32$ and
$1\leq a\leq \log_2 m$.}
\label{tab:gold}
\end{table}

Most of these forms have a nontrivial radical. The statement is therefore not
``nondegenerate Gold form'' but ``positive-rank quadratic form with a
nonsingular core whose Arf value is computed after symplectic reduction.''

\subsection{A broader trace family}

The coefficient-$1$ Gold form is only one trace monomial. A standard trace
description of the quadratic part of Boolean functions over $\F_{2^m}$ uses
terms
\[
  \Tr(c_i x^{1+2^i})
\]
plus affine terms and, in even degree, a middle half-trace term. The monomials
displayed here are still game-operation-built in the nimber backend:
$c_i x^{1+2^i}=c_i\nimmul x\nimmul x^{2^i}$, followed by trace/XOR.

The current implementation does not claim to package every Boolean quadratic
function. The probe \texttt{experiments/gold\_family\_survey.py} deliberately
omits the middle term and affine bookkeeping. What it does check, exhaustively
over $\F_{256}$, is that scaled Gold components
\[
  Q_{\lambda,a}(x)=\Tr(\lambda x^{1+2^a})
\]
already reach nondegenerate (bent) forms for many $\lambda$. For APN Gold
exponents in that field, the observed count matches the classical
$2(2^m-1)/3$ bent-component count. This matters for the open game question
because a bent form has trivial radical: the loopy radical route collapses to
$\{0\}$, and the frame-blind symplectic no-go applies without a degenerate
layer to hide in.

\begin{remark}[a vacuous small target]\label{rem:vacuous}
At $m=4$, $a=1$ the zero set $\{Q_1=0\}\subset\F_{16}$ is the subfield
$\F_4$, an $\F_2$-\emph{subspace} of dimension $2$: the form is genuinely
quadratic (polar rank $2$), but its zero set is affine, so as a target for
``is this $P$-set a genuine quadric'' it is vacuous -- a linear rule realizes
it. The smallest honest unscaled instance is $(8,1)$. The repository's
quadric-fitting layer (\texttt{src/forms/quadric\_fit.rs}) reports exactly
this distinction (genuine polar rank versus affine flat), and every probe
below that consumes a fitted quadric is read with that guard in mind.
\end{remark}

\section{Arf as conditional win-bias}\label{sec:bias}

The relevant counting theorem is classical. If $Q$ is nonsingular on
$\F_2^{2r}$, then
\[
  \#\{x:Q(x)=0\}
  =
  2^{2r-1}+(-1)^{\Arf(Q)}2^{r-1}.
\]
For degenerate forms, the code uses the standard radical-adjusted count: an
anisotropic radical balances the values, while an isotropic radical scales the
bias of the nonsingular core.

The game interpretation is conditional:

\begin{quote}
If a game has $P$-positions exactly $\{x:Q(x)=0\}$, then the Arf invariant is
the sign of the bias between second-player wins and first-player wins.
\end{quote}

This is not a claim that such a game has been found. It is a target for a game
semantics.

\begin{table}[h]
\centering
\begin{tabular}{rrcccr}
\toprule
field & $a$ & $\Arf$ & rank & $\#\{Q=0\}$ & bias \\
\midrule
$\F_{2^{4}}$  & $1$ & $1$ & $2$  & $4$     & $-4$ \\
$\F_{2^{8}}$  & $1$ & $1$ & $6$  & $112$   & $-16$ \\
$\F_{2^{8}}$  & $2$ & $1$ & $4$  & $96$    & $-32$ \\
$\F_{2^{16}}$ & $1$ & $1$ & $14$ & $32512$ & $-256$ \\
$\F_{2^{16}}$ & $4$ & $1$ & $8$  & $30720$ & $-2048$ \\
\bottomrule
\end{tabular}
\caption{Representative zero-counts from
\texttt{experiments/arf\_win\_bias.py}. The bias is
$\#\{Q=0\}-2^{m-1}$.}
\label{tab:bias}
\end{table}

\section{The linear ceiling}\label{sec:linear}

For a normal-play disjunctive sum of impartial games, the P-condition is
linear in Grundy coordinates:
\[
  g_1\oplus\cdots\oplus g_n=0.
\]
So normal-play sums give linear P-sets, not genuine quadrics. The
coin-turning architecture -- the very architecture whose product theorem
defines the nim product -- is subject to the same ceiling.

\begin{nogo}[coin-turning linearity; standard math, machine cross-checked]
\label{thm:A}
Every Winning Ways coin-turning game has configuration $P$-set equal to the
kernel of an $\F_2$-linear nimber-valued map. In particular no coin-turning
$P$-set is a quadric of positive polar rank.
\end{nogo}

\begin{proof}[Proof sketch]
A coin-turning position is a finite set of heads; by the Winning Ways
coin-turning theory~\cite{WW}, the Grundy value of a configuration is the XOR
of the single-coin Grundy values $g(i)$ of its heads. The $P$-set is therefore
$\{x:\bigoplus_{i\in x}g(i)=0\}$, the kernel of the $\F_2$-linear map
$x\mapsto\bigoplus_i x_i\,g(i)$. The repository cross-check
(\texttt{experiments/gold/octal\_attack.py}) verifies XOR-additivity by mex
recursion, exhaustively over all $32{,}768$ companion-family choices on $4$
coins.
\end{proof}

Theorem~\ref{thm:A} has a sharp classical complement: the linear ceiling is
exactly where coding theory already harvests games.

\begin{remark}[the lexicode shadow; standard math and interpretation]
\label{rem:lexicode}
Conway--Sloane lexicodes~\cite{ConwaySloane86} are built by the greedy
lexicographic rule -- the mex rule -- and the 1986 paper realizes them as the
Grundy-value-$0$ positions of associated turning-game move structures; binary
lexicodes are linear \emph{because of} Sprague--Grundy theory (XOR-closure is
a game theorem there, not a coding theorem). They include the Hamming codes
and the length-$24$, $d=8$ extended binary Golay code; over base $2^k$ they
are closed under nim-addition, and they are linear at the Fermat-power bases
$2^{2^k}$ -- exactly the sizes at which nim-multiplication makes the ordinals
below the base a field. Read together with Theorem~\ref{thm:A}: lexicodes are
the \emph{floor} (rich linear codes naturally realized as Grundy-$0$ sets) and
Theorem~\ref{thm:A} is the \emph{ceiling} (coin-turning realizes only linear
$P$-sets); floor and ceiling meet at linear. The open problem of this note is
whether the lexicode phenomenon admits a quadratic refinement, and the
obstruction is precisely the XOR-closure failure that the polar form
measures.
\end{remark}

For a quadratic form in characteristic $2$,
\[
  Q(u+v)=Q(u)+Q(v)+B(u,v).
\]
Thus for $u,v\in\{Q=0\}$,
\[
  u+v\in\{Q=0\}\quad\Longleftrightarrow\quad B(u,v)=0.
\]
The polar form is exactly the obstruction to XOR-closure. The first two layers
are game-built: XOR is disjunctive sum and $B$ is built from nim products. The
missing layer is a play rule that integrates that bilinear data into the
quadratic zero set.

\section{Routes and obstructions}\label{sec:blocks}

\subsection{Test benches}\label{sec:benches}

The repo contains several probes. They are useful for narrowing the target, but
they are not a proof that no natural game exists.

\begin{itemize}[leftmargin=*]
\item \texttt{fit\_f2\_quadratic} tests whether a subset of $\F_2^k$ is the zero
set of a quadratic polynomial, and whether the result has nonzero polar rank
(distinguishing genuine quadrics from affine flats; cf.\
Remark~\ref{rem:vacuous}).
\item A bounded misere quotient computer tests small quotient P-sets. The octal
sweep checked $292$ octal codes, heap cutoff at most $4$, and found quotient
orders $2,6,10,12,14$, with no elementary abelian $2$-group of rank at least
$2$. Theorem~\ref{thm:C} below shows this emptiness is forced, not a sampling
accident.
\item \texttt{experiments/misere\_kernel.py} checks the kernel obstruction on
the smallest wild misere quotient $R_8$: the canonical group/kernel shadow is
linear, while the genuinely misere P-element lies outside that vector-space
part. Corollary~\ref{cor:kernel} below shows the linearity is forced by
commutativity.
\item An interactive-kernel solver tests finite move graphs. It confirms that
an ad hoc acyclic game can realize any subset, and that a rule referencing $Q$
itself realizes $\{Q=0\}$ tautologically. The tested $B$-coupled rules do not
produce the Gold zero set.
\item \texttt{examples/loopy\_quadric.rs} tests cyclic move graphs with Draw
sets. The symmetric $B$-coupled rule has Loss-set equal to the radical $R(B)$,
so it reproduces the obstruction rather than the Gold zero set; this is the
special case of Theorem~\ref{thm:G} below, and its small $(4,1)$ ``hit'' is
the radical coincidence $|\{Q_1=0\}|=4=|R(B)|$ there.
\item \texttt{examples/bent\_route.rs} repeats route probes on a bent scaled
Gold component. In that example, a $B$ plus coordinate-frame rule reaches a
quadric of the correct Arf class but not the specific Gold zero set; the naive
local-field/spin-flip completion leaves the quadric variety.
\end{itemize}

Except for the frame-blind proposition below and the lettered theorems of
Section~\ref{sec:nogos}, these are bounded or example-level negative results.

\subsection{Frame-blind no-go}

There is a clean obstruction for rules that see only a nondegenerate
symplectic form and no frame.

\begin{proposition}\label{prop:nogo}
Let $B$ be a nondegenerate alternating form on $V=\F_2^{2r}$ with $r\geq 2$.
If a finite normal-play game has position set $V$ and a move relation invariant
under $\Sp(B)$, then its $P$-set is not a nondegenerate quadric.
\end{proposition}

\begin{proof}
Each element of $\Sp(B)$ is an automorphism of the move graph, so
the retrograde outcome function is $\Sp(B)$-invariant. The P-set
is therefore a union of symplectic orbits. For $r\geq2$, the symplectic group is
transitive on $V\setminus\{0\}$~\cite{Taylor}, so the invariant subsets are only
$\varnothing$, $\{0\}$, $V\setminus\{0\}$, and $V$. A nondegenerate quadric in
dimension at least $4$ has $2^{2r-1}\pm2^{r-1}$ points, which is none of these.
\end{proof}

The hypothesis matters. Coordinate-dependent rules have already broken
$\Sp(B)$ symmetry, so this proposition does not explain every
negative probe. It explains only genuinely frame-blind $B$-only rules.
Section~\ref{sec:nogos} extends the obstruction far beyond this symmetric
class.

\subsection{The diagonal framing}\label{sec:framing}

Once a coordinate frame is allowed, a canonical-looking quadratic refinement of
$B$ is
\[
  Q_{\mathrm{frame}}(v)=\sum_{i<j}B(e_i,e_j)v_i v_j.
\]
For the Gold polar forms tested in this repo, the Gold form decomposes as
\[
  Q_a(v)=Q_{\mathrm{frame}}(v)+\ell_{\mathrm{diag}}(v),
  \qquad
  \ell_{\mathrm{diag}}(v)=\sum_i Q_a(e_i)v_i.
\]
In the tested genuinely quadratic Gold cases up to $\F_{2^{16}}$,
$Q_{\mathrm{frame}}$ has Arf $0$ and the diagonal term changes the form to the
Gold Arf-$1$ class.

\begin{remark}
This is an experimental pattern for the Gold polar forms tested here, not a
theorem about every alternating form plus every coordinate frame. Random
nondegenerate alternating forms can give a zero-diagonal frame quadric of Arf
$1$.
\end{remark}

The diagonal $q_i=Q_a(e_i)$ is therefore the named missing datum. Two results
of the parallel run locate it more precisely.

\begin{proposition}[game-native diagonal source for $a=1$; implemented and
tested]\label{prop:wp}
Let $m=2^k$ and let $w\in\F_{2^m}$ be the XOR of the Fermat coins
$2^{2^t}$ for $t=1,\dots,k-1$. Then the $a=1$ Gold diagonal satisfies
\[
  q_i^{(m,1)}=Q_1(e_i)=\Tr\!\bigl(\wp(w)\,e_i\bigr),
  \qquad
  \wp(w)=w\nimmul w\oplus w,
\]
i.e.\ the trace-dual element of the diagonal is $\wp(w)$. Verified at
$m=4,8,16,32$ via the trace-pairing uniqueness of the dual element
(\texttt{experiments/gold/gold\_diag\_probe.py}).
\end{proposition}

This gives a game-native, tower-uniform, commutative source for the $a=1$
diagonal: one nim-squaring, one XOR, one trace -- all inside the licensed
operation chain of Section~\ref{sec:gold}. For even $a$ the dual
$\lambda_a^{(m)}$ provably drifts at each tower doubling
($\lambda_2$: $0,6,102,24582$; $\lambda_4$: $31,8030$), and no
$\wp$-preimage family has been named beyond $m=8$; the even-$a$ analogue is
open.

\begin{remark}[subfield vanishing; standard math]
$q_i^{(m,a)}=0$ for all $i<m/2$: the Fermat basis elements at lower indices,
their Frobenius iterates, and their nim products stay inside the half-field,
whose absolute trace from an even-degree extension vanishes. The Gold diagonal
lives on the top Fermat layer only.
\end{remark}

\begin{question}
Is the diagonal framing
\[
  q_i=Q_a(e_i)=\Tr(e_i^{1+2^a})
\]
itself game-natural for all $a$, in a way that yields a natural game with
P-set $\{x:Q_a(x)=0\}$? Proposition~\ref{prop:wp} answers the
\emph{sourcing} half affirmatively for $a=1$; whether any acceptable source
yields a non-tautological \emph{play rule} is the subject of the remaining
sections.
\end{question}

\section{The extraspecial reframing}\label{sec:extraspecial}

The diagonal-framing question can be restated in group-theoretic terms, and
the restatement sharpens it. The dictionary is classical -- the
characteristic-$2$ Heisenberg/extraspecial picture of quadratic
forms~\cite{Quillen71,Gorenstein80} -- and this section proves the pieces we
use, to keep the note self-contained. Claim levels: Lemmas~\ref{lem:extdict},
\ref{lem:arftypes}, \ref{lem:abelian} and
Proposition~\ref{prop:between} are standard mathematics; the $R_8$
instantiation in Corollary~\ref{cor:kernel} is implemented and tested
(\path{experiments/misere_kernel.py}); reading $E$-equivariance as the right
naturality criterion is interpretation; the existence of a game-native model
of the extension is open.

Throughout, $V=\F_2^n$ written additively, and $Z=\{1,z\}\cong\mathbb{Z}/2$
written multiplicatively. A \emph{quadratic map} $Q\colon V\to\F_2$ is a
function with $Q(0)=0$ whose polarization
$B(x,y)=Q(x+y)+Q(x)+Q(y)$ is bilinear; $B$ is then alternating (hence
symmetric), and $Q$ may have any diagonal $q_i=Q(e_i)$, exactly the
char-$2$ datum the engine keeps separate from $b$.

\subsection{Quadratic forms are central extensions}

\begin{lemma}[extraspecial dictionary]\label{lem:extdict}
Let $1\to Z\to E\xrightarrow{\ \pi\ }V\to 1$ be a central extension of the
elementary abelian $2$-group $V$ by $Z\cong\mathbb{Z}/2$.
\begin{enumerate}[label=(\roman*),leftmargin=*]
\item For $x\in V$ and any lift $\tilde x\in\pi^{-1}(x)$ one has
$\tilde x^2\in Z$, and $\tilde x^2$ is independent of the lift; the
\emph{squaring form} $Q_E\colon V\to\F_2$ defined by
$\tilde x^2=z^{Q_E(x)}$ is well defined, as is the \emph{commutator pairing}
$B_E$ defined by $[\tilde x,\tilde y]=z^{B_E(x,y)}$, which is bilinear and
alternating. They satisfy
\[
(\tilde x\tilde y)^2=\tilde x^2\,\tilde y^2\,[\tilde x,\tilde y],
\qquad\text{equivalently}\qquad
Q_E(x+y)=Q_E(x)+Q_E(y)+B_E(x,y),
\]
so $Q_E$ is a quadratic map with polar form $B_E$.
\item Conversely, every quadratic map $Q$ arises: fixing a basis
$e_1,\dots,e_n$ of $V$, the bilinear $2$-cocycle
\[
\varphi(x,y)=\sum_i x_iy_i\,Q(e_i)+\sum_{i>j}x_iy_j\,B(e_i,e_j)
\]
defines a group $E_Q=Z\times V$ with
$(s,x)(t,y)=(s+t+\varphi(x,y),\,x+y)$, whose squaring form is $Q$ and whose
commutator pairing is the polar form $B$.
\item The assignment $Q\mapsto[E_Q]$ is a bijection from quadratic maps on
$V$ to equivalence classes of central extensions of $V$ by $\mathbb{Z}/2$.
\item $E_Q$ is abelian iff $B=0$ (iff $Q$ is linear); and for $n\geq 1$,
$E_Q$ is \emph{extraspecial} (i.e.\ $Z(E)=[E,E]=\Phi(E)\cong\mathbb{Z}/2$)
iff $B$ is nondegenerate, which forces $n=2r$ even and $|E_Q|=2^{1+2r}$.
\end{enumerate}
\end{lemma}

\begin{proof}
(i) $\pi(\tilde x^2)=2x=0$, so $\tilde x^2\in Z$; replacing $\tilde x$ by
$z\tilde x$ multiplies the square by $z^2=1$. Commutators of lifts lie in
$Z$ because $V$ is abelian, are central, and are unchanged by central
retagging of the lifts; centrality of commutators gives bimultiplicativity,
and $[\tilde x,\tilde x]=1$ gives $B_E(x,x)=0$. For the identity: with
$c=[\tilde y,\tilde x]$ central,
$\tilde x\tilde y\tilde x\tilde y
=\tilde x(\tilde x\tilde y\,c)\tilde y=\tilde x^2\tilde y^2c$,
and $c=c^{-1}=[\tilde x,\tilde y]$ since $Z$ has exponent $2$. As
$\tilde x\tilde y$ lifts $x+y$, the additive identity follows.

(ii) A bilinear map is a $2$-cocycle, so $E_Q$ is a group with center
containing $Z=\{(0,0),(1,0)\}$. Squares: $(s,x)^2=(\varphi(x,x),0)$ and
$\varphi(x,x)=\sum_i x_iQ(e_i)+\sum_{i>j}x_ix_jB(e_i,e_j)=Q(x)$ by the
polarization expansion of $Q$ in the basis. Commutators:
$(s,x)(t,y)(s,x)^{-1}(t,y)^{-1}=(\varphi(x,y)+\varphi(y,x),0)
=(B(x,y),0)$, using that the symmetrization of $\varphi$ is $B$ (the
diagonal terms cancel and $B$ is symmetric).

(iii) Equivalent extensions have cohomologous cocycles, and a coboundary
$\delta\psi(x,y)=\psi(x)+\psi(y)+\psi(x+y)$ has zero diagonal
($\delta\psi(x,x)=\psi(0)=0$), so the squaring form is an invariant of the
class; (ii) gives surjectivity. Injectivity: if two extensions have equal
squaring form $Q$, choose sections and cocycles $\varphi,\varphi'$; both
have diagonal $Q$, and both have symmetrization equal to the polar form of
$Q$, so $d=\varphi+\varphi'$ is a symmetric cocycle with zero diagonal. The
extension $E_d$ is then abelian of exponent $2$, i.e.\ an $\F_2$-vector
space, so $1\to Z\to E_d\to V\to 1$ splits and $d$ is a coboundary; the two
extensions are equivalent.

(iv) Commutators generate $z^{\,\mathrm{im}\,B}$, so $E_Q$ is abelian iff
$B=0$. In general $Z(E_Q)=\pi^{-1}(\rad B)$, since $\tilde x$
is central iff $B(x,\cdot)=0$; thus $Z(E_Q)=Z$ iff $B$ is nondegenerate. If
so, $B\neq0$ gives $[E,E]=Z$, and $\Phi(E)=E^2[E,E]=Z$ because all squares
lie in $Z$ and $E/Z$ is elementary abelian. A nondegenerate alternating
form exists only in even dimension.
\end{proof}

\begin{remark}
This is the characteristic-$2$ Heisenberg picture: $E_Q$ is the Heisenberg
group of $(V,B)$ and $Q$ selects the central extension; cohomologically,
(iii) is the standard identification of $H^2(V;\F_2)$ with quadratic maps
$V\to\F_2$~\cite{Quillen71}. The same finite-quadratic-module data drives
the Weil $S$/$T$ matrices in the library's integral layer. Note that the
dictionary is exactly the engine's discipline in group clothing: the
diagonal $q_i$ (squares) and the off-diagonal $b_{ij}$ (commutators) are
independent central data, and collapsing them is collapsing $E_Q$ onto an
abelian quotient.
\end{remark}

\subsection{Arf classifies the two extraspecial types}

Write $H$ for the hyperbolic plane ($Q(e_1)=Q(e_2)=0$, $B(e_1,e_2)=1$;
$\Arf=0$) and $A$ for the anisotropic plane ($Q\equiv1$ on $V\setminus\{0\}$;
$\Arf=1$), and $\circ$ for the central product (identify the centers).

\begin{lemma}[Arf = the two types]\label{lem:arftypes}
Let $B$ be nondegenerate, $\dim V=2r\geq2$.
\begin{enumerate}[label=(\roman*),leftmargin=*]
\item $E_{Q\perp Q'}\cong E_Q\circ E_{Q'}$.
\item $E_H\cong D_8$ and $E_A\cong Q_8$.
\item Every extraspecial $2$-group of order $2^{1+2r}$ is isomorphic to
$E_Q$ for a nondegenerate $Q$, and
$E_Q\cong E_{Q'}$ iff $(V,Q)\cong(V',Q')$ iff the ranks and Arf invariants
agree~\cite{Dickson01,Arf41}. Hence there are exactly two extraspecial
groups of order $2^{1+2r}$:
\[
E^+_{2^{1+2r}}\;=\;D_8^{\circ r}\ \ (\Arf=0,\ Q\cong H^{\perp r}),
\qquad
E^-_{2^{1+2r}}\;=\;D_8^{\circ(r-1)}\circ Q_8\ \ (\Arf=1,\ Q\cong
H^{\perp(r-1)}\perp A).
\]
In particular $D_8\circ D_8\cong Q_8\circ Q_8$, the group avatar of
$H\perp H\cong A\perp A$ (additivity of $\Arf$).
\item (Census.) With $\varepsilon=(-1)^{\Arf Q}$,
\[
\#\{g\in E_Q: g^2=1\}=2\,\#\{x:Q(x)=0\}=2^{2r}+\varepsilon\,2^{r},
\]
so the element-order census of $E_Q$ is the zero-count bias of
Section~\ref{sec:bias} read multiplicatively, and it distinguishes the two
types.
\end{enumerate}
\end{lemma}

\begin{proof}
(i) The cocycle $\varphi\oplus\varphi'$ on $V\oplus V'$ is bilinear with
diagonal $Q\perp Q'$, and
$(s,(x,x'))\mapsto[((s,x),(0,x'))]$ is an isomorphism onto
$(E_Q\times E_{Q'})/\langle(z,z')\rangle$.

(ii) $E_H$ has order $8$, is nonabelian ($B\neq0$), and has a noncentral
involution (any lift of $e_1$ squares to $z^{Q(e_1)}=1$), so $E_H\cong D_8$;
in $E_A$ every noncentral element squares to $z$ (as $Q\equiv1$ off $0$), so
$z$ is the unique involution and $E_A\cong Q_8$.

(iii) For $E$ extraspecial, $V=E/Z(E)$ is elementary abelian (as
$\Phi(E)=Z(E)$), so Lemma~\ref{lem:extdict} applies to
$1\to Z(E)\to E\to V\to1$ and gives $E\cong E_{Q_E}$ with $B_{Q_E}$
nondegenerate ($Z(E)$ exactly central). If $\psi\colon E_Q\to E_{Q'}$ is an
isomorphism, it preserves the (characteristic) centers, fixes $z\mapsto z'$
(the unique nontrivial central element), and induces an $\F_2$-linear map
$\bar\psi\colon V\to V'$ with
$Q'(\bar\psi x)$ read from $\psi(\tilde x)^2=\psi(\tilde x^2)$ -- an
isometry. Conversely an isometry $g$ transports the cocycle:
$\varphi'\circ(g\times g)$ has diagonal $Q'\circ g=Q$, hence by
Lemma~\ref{lem:extdict}(iii) yields an extension equivalent to $E_Q$, and
$(s,x)\mapsto(s,gx)$ closes the isomorphism. Dickson's classification of
nonsingular quadratic forms over $\F_2$ (two isometry classes per even
rank, separated by $\Arf$) finishes; the central-product normal forms
follow from (i), (ii) and Witt decomposition.

(iv) The two lifts of $x$ both square to $z^{Q(x)}$, so they are
involutions or the identity iff $Q(x)=0$; now apply the zero-count formula
of Section~\ref{sec:bias}. The counts differ for $\varepsilon=\pm1$, so
$E^+\not\cong E^-$.
\end{proof}

\subsection{The abelian obstruction}

\begin{lemma}[abelian obstruction]\label{lem:abelian}
Let $M$ be a commutative monoid, $Z=\{1,z\}\subseteq M$ a two-element
subgroup, $N\subseteq M$ a submonoid containing $Z$, and
$\pi\colon N\twoheadrightarrow V$ a surjective monoid homomorphism onto an
$\F_2$-vector space with $\pi^{-1}(\pi(m))=mZ$ for all $m\in N$. Then the
squaring form $Q\colon V\to\F_2$, defined by $\tilde m^2=z^{Q(x)}$ for any
$\tilde m\in\pi^{-1}(x)$, is well defined and \emph{$\F_2$-linear}; its
polar form vanishes identically. Consequently no commutative monoid
realizes, through its own squaring map, a quadratic form with nonzero --
a fortiori nondegenerate -- characteristic-$2$ polar form. Equivalently, by
Lemma~\ref{lem:extdict}(iv): if $B\neq0$, the extension $E_Q$ is
nonabelian and admits no model $(N,Z,\pi)$ inside any commutative monoid.
\end{lemma}

\begin{proof}
$\pi(m^2)=2\pi(m)=0$, so $m^2\in\pi^{-1}(0)=Z$, and
$(zm)^2=z^2m^2=m^2$, so $Q$ is well defined. Commutativity gives
$(mn)^2=m^2n^2$; since $\tilde m\tilde n$ is a preimage of $x+y$, this is
$Q(x+y)=Q(x)+Q(y)$, i.e.\ $B\equiv0$.
\end{proof}

\begin{corollary}[misère quotients: the kernel shadow is linear]\label{cor:kernel}
Let $\mathcal{Q}$ be a finite misère quotient -- a finite commutative
monoid -- with kernel $K$, the maximal subgroup of $\mathcal{Q}$ (the
mutual-divisibility class of the product of all idempotents)~\cite{PS}.
\begin{enumerate}[label=(\roman*),leftmargin=*]
\item Every configuration $(N,Z,\pi)$ as in Lemma~\ref{lem:abelian} inside
$\mathcal{Q}$ has linear squaring form: the intrinsic multiplication of a
misère quotient cannot supply a quadratic refinement with $B\neq0$, on $K$
or anywhere else in $\mathcal{Q}$.
\item Under the regularity hypothesis of \cite[Thm.~6.4]{PS} (satisfied by
the regular finite quotients arising in practice), $K\cong(\mathbb{Z}/2)^k$
is the group of normal-play Grundy values and the $P$-portion meets $K$ in
the XOR-linear normal-play set.
\item On the smallest wild quotient
$R_8=\langle a,b,c\mid a^2=1,\ b^3=b,\ bc=ab,\ c^2=b^2\rangle$ with
$P=\{a,b^2\}$: the idempotents are $\{1,b^2\}$, the kernel is
$K=\{b^2,b,ab,ab^2\}\cong(\mathbb{Z}/2)^2$ with identity $b^2$,
$P\cap K=\{b^2\}\mapsto\{0\}$ is linear, and the genuinely misère
$P$-element $a$ lies outside $K$. (Implemented and tested:
\path{experiments/misere_kernel.py}.)
\end{enumerate}
So the linear obstruction observed on $R_8$ is forced by commutativity, not
an accident of the example: a nondegenerate polar form is the commutator
pairing of a nonabelian group, and a misère quotient has none to offer.
\end{corollary}

\begin{proof}
(i) is Lemma~\ref{lem:abelian} applied to submonoids of $\mathcal{Q}$,
which are commutative. (ii) is quoted from \cite{PS}. (iii) is a finite
verification.
\end{proof}

\subsection{The integral shadow}\label{sec:integral}

The same extension data has an integral-lattice avatar, already shipped in
the library (interpretation, verified through the
\texttt{DiscriminantForm} + \texttt{arf\_nimber} pipeline): the Gold
nonsingular core of rank $2r$ and Arf $\varepsilon$ is the discriminant form
of the even lattice $U(2)^{r}$ ($\varepsilon=0$) or
$U(2)^{r-1}\oplus D_4$ ($\varepsilon=1$); the Milgram phase is
$4\cdot\Arf\bmod 8$; the Weil $S$-matrix prefactor is $(-1)^{\Arf}$; and the
Frobenius--Schur indicator of the Heisenberg representation of $E_Q$ is
$(-1)^{\Arf}$. The metaplectic relation $(ST)^3=S^2$ holds iff the diagonal
$T$ uses a quadratic refinement of the correct Arf class -- it selects the
Arf \emph{class}, not the member. So the integral bridge carries exactly one
game-relevant bit; it cannot, by itself, distinguish the Gold form from any
other refinement in its class.

\subsection{A Tier-2 naturality screen: $E$-equivariance}

Proposition~\ref{prop:nogo} kills rules that are blind to everything but
the symplectic structure (Tier~1), while per-position evaluator circuits
for $Q_a$ realize the quadric tautologically (Tier~3). The extraspecial
dictionary suggests a criterion strictly between the two: the rule should
see the extension $E_Q$ -- which carries the diagonal data $q_i$
structurally, as squares -- but only up to its automorphisms, so that no
basis, field structure, or evaluation circuit can be smuggled in.

\begin{definition}[uniform rules; the three tiers]\label{def:tier2}
Let $Q$ be nondegenerate on $V=\F_2^{2r}$ with polar form $B$, and
$E=E_Q$, $\pi\colon E\to V$, $\Sigma=\{x\in V: Q(x)=0\}$.
\begin{enumerate}[label=(\alph*),leftmargin=*]
\item A \emph{uniform rule on a finite set $X$} is a single move relation
$M\subseteq X\times X$, fixed once for all positions; for normal play we
require the move digraph acyclic, and outcomes are computed by the usual
retrograde recursion. (For loopy or misère readings, replace the outcome
recursion; the equivariance constraint below applies verbatim, since
outcome classes are invariants of digraph automorphisms.)
\item A uniform rule on $E$ is \emph{$E$-equivariant} if every
$\alpha\in\Aut(E)$ is an automorphism of the move digraph.
Its \emph{shadow} is $\pi(P)\subseteq V$, where $P$ is its $P$-set; it
\emph{realizes} $T\subseteq V$ if $\pi(P)=T$.
\item \emph{Tier 1} ($\Sp(B)$-blind): a uniform rule on $V$
invariant under $\Sp(B)$, as in
Proposition~\ref{prop:nogo}. \emph{Tier 3} ($Q_a$-evaluation): a family of
games $\{\Gamma_x\}_{x\in V}$ whose structure varies with the designated
input $x$ (e.g.\ the trace-circuit evaluator); this is not a uniform rule.
\emph{Tier 2 screen}: a route to the Gold quadric is Tier-2 natural
\emph{only if} its positions and moves lift to an $E$-equivariant uniform
rule on $E_{Q_a}$ -- with $Q_a$ restricted to its nonsingular core when the
form is degenerate, matching the classifier's radical discipline -- that
realizes $\{Q_a=0\}$.
\end{enumerate}
\end{definition}

\begin{proposition}[the screen sits strictly between the solved tiers]\label{prop:between}
Let $Q$ be nondegenerate on $V=\F_2^{2r}$, $r\geq2$, and $E=E_Q$.
\begin{enumerate}[label=(\roman*),leftmargin=*]
\item The image of $\Aut(E)\to \GL(V)$ is exactly the
orthogonal group $\Orth(Q)$, with kernel the central twists
$g\mapsto g\,z^{\ell(\pi g)}$, $\ell\in V^{*}$, which equal
$\Inn(E)\cong V$ because $B$ is nondegenerate; thus
$\Out(E)\cong \Orth(Q)$ (cf.~\cite{Winter72}). The
$\Aut(E)$-orbits on $E$ are
\[
\{1\},\qquad \{z\},\qquad \pi^{-1}(\Sigma\setminus\{0\}),\qquad
\pi^{-1}(V\setminus\Sigma).
\]
\item Hence the $P$-set of any $E$-equivariant uniform rule is a union of
these four orbits, and its shadow is one of the eight unions of $\{0\}$,
$\Sigma\setminus\{0\}$, $V\setminus\Sigma$. The quadric $\Sigma$ is among
them: the Tier-1 exclusion does not apply.
\item (Tier 1 $\subsetneq$ Tier 2.) Every $\Sp(B)$-invariant
uniform rule on $V$ pulls back along $\pi$ to an $E$-equivariant rule with
$P$-set $\pi^{-1}$ of the original, whose shadow is therefore one of
$\varnothing$, $\{0\}$, $V\setminus\{0\}$, $V$ -- never the quadric
(Proposition~\ref{prop:nogo} in this language). The containment is strict:
the \emph{squaring rule}
\[
g\to h\ \text{legal}\iff g^2=z\ \text{and}\ h^2=1
\]
(``move from any order-$4$ position to any position of order at most
$2$'') is $E$-equivariant and acyclic, and its $P$-set is the involution
locus $\{g:g^2=1\}=\pi^{-1}(\Sigma)$, with shadow exactly $\Sigma$.
\item (Tier 2 $\subsetneq$ Tier 3.) Up to the isometry induced by any group
isomorphism, the family of $E$-equivariant rules and their shadows depends
only on $(r,\Arf Q)$ (Lemma~\ref{lem:arftypes}(iii)). In particular the
screen cannot separate Gold forms of equal core rank and Arf invariant, and
an $E$-equivariant rule has no access to $m$, the exponent $a$, the field
multiplication, the coordinate frame $q_i=Q_a(e_i)$, or any evaluation
circuit. Conversely, arbitrary subsets of $V$ -- all realizable by ad hoc
acyclic lookup games and by Tier-3 evaluators
(Section~\ref{sec:benches}) -- are not shadows of $E$-equivariant rules once
they leave the eight-set list of (ii).
\end{enumerate}
\end{proposition}

\begin{proof}
(i) Automorphisms preserve $Z=Z(E)$ and fix $z$
($\Aut(\mathbb{Z}/2)=1$), so the induced map preserves the
squaring form: $Q(\bar\alpha x)$ is read from
$\alpha(\tilde x)^2=\alpha(\tilde x^2)=z^{Q(x)}$; the image lies in
$\Orth(Q)$. Surjectivity: an isometry $g$ transports the standard cocycle as in
the proof of Lemma~\ref{lem:arftypes}(iii), producing an automorphism of
$E_Q$ inducing $g$. The kernel consists of maps $g\mapsto g\,z^{\ell(\pi g)}$
with $\ell$ linear (multiplicativity forces additivity of $\ell$); inner
automorphisms are exactly the twists $\ell=B(\bar g,\cdot)$, and
nondegeneracy makes $\bar g\mapsto B(\bar g,\cdot)$ onto $V^{*}$. Orbits:
$\Orth(Q)$ is transitive on $\Sigma\setminus\{0\}$ and on $V\setminus\Sigma$ by
Witt's extension theorem in its characteristic-$2$ quadratic-space form
\cite{Taylor} (both sets are nonempty for $r\geq2$); the two lifts of any
$x\neq0$ are conjugate under conjugation by $\tilde g$ with $B(\bar g,x)=1$;
and $1$, $z$ are fixed by everything.

(ii) Digraph automorphisms preserve the retrograde outcome recursion, so
$P$ is $\Aut(E)$-invariant; apply (i) and project.

(iii) The pullback $M=\{(g,h):(\pi g,\pi h)\in R\}$ of an
$\Sp(B)$-invariant rule $R$ is preserved by
$\Aut(E)$, since the induced action on $V$ lies in
$\Orth(Q)\subseteq\Sp(B)$; it is acyclic when $R$ is, and an easy
induction along the acyclic rank gives
$P(M)=\pi^{-1}(P(R))$. By the transitivity of
$\Sp(B)$ on $V\setminus\{0\}$ for $r\geq2$, $P(R)$ is a union
of $\{0\}$ and $V\setminus\{0\}$. For the squaring rule: automorphisms
preserve element orders, so the rule is $E$-equivariant; positions with
$g^2=1$ are terminal, hence $P$; positions with $g^2=z$ have a move (to
$1$), hence $N$. Its shadow $\Sigma$ satisfies
$1<|\Sigma|=2^{2r-1}+\varepsilon2^{r-1}<2^{2r}-1$ for $r\geq2$, so it is
none of the four Tier-1 shadows.

(iv) If $\psi\colon E_Q\to E_{Q'}$ is an isomorphism (it exists iff ranks
and Arf agree, Lemma~\ref{lem:arftypes}(iii)), then
$M\mapsto(\psi\times\psi)(M)$ is a bijection between $E$-equivariant rules
carrying $P$-sets to $P$-sets and shadows to shadows through the induced
isometry $\bar\psi$; in particular it carries rules realizing $\{Q=0\}$ to
rules realizing $\{Q'=0\}$. The shadow constraint of (ii) excludes all but
eight subsets, while Tier-3 constructions realize any subset; the
containment is proper.
\end{proof}

\begin{remark}[what the screen does and does not settle]\label{rem:screen}
Equivariance is a screen, not a construction. Over the \emph{abstract}
group $E$ the screen is satisfiable -- Proposition~\ref{prop:between}(iii)
exhibits the squaring rule, which is nothing but the extraspecial squaring
map of Lemma~\ref{lem:extdict} read as a one-move relation. That rule
consults only the multiplication of $E$, so it is a $Q$-evaluator exactly
insofar as $E$ itself is fed in by hand. The reframing therefore relocates
the open problem: instead of \emph{find a play rule that computes $Q_a$},
it asks \emph{exhibit a game-native model of the extension $E_{Q_a}$} --
positions built from game values, multiplication from game constructions --
after which the rule layer is canonical. Lemma~\ref{lem:abelian} is the
sharp constraint on that search: no commutative value-world, in particular
no misère-quotient kernel (Corollary~\ref{cor:kernel}), can host $E$ once
$B\neq0$. The noncommutativity must enter from a structurally available
asymmetry, and the one that normal, misère, and partizan play all possess
-- and that the symmetric polar form $B$ discards -- is the
first-/second-player asymmetry of the move relation. Claim
levels: the proposition is standard mathematics; treating
$E$-equivariance as \emph{the} naturality criterion is interpretation;
the existence of a game-native $E$ is open. (For $r=1$ the screen degenerates:
on the anisotropic plane $\Sigma=\{0\}$ and $\Orth(Q)=\Sp(B)$,
so the statement is vacuous there; the Gold targets of interest have
$r\geq2$ cores.) Theorem~\ref{thm:E} below sharpens the screen further:
read as \emph{invariance} it admits no live middle at all, so the
$E$-structure must enter \emph{covariantly}, in the dynamics.
\end{remark}

\begin{question}
Does any game-native construction realize the extraspecial extension
$E_{Q_a}$ -- equivalently, produce the diagonal data $q_i=Q_a(e_i)$ as
squares in a noncommutative game-built structure -- without an evaluation
circuit for $Q_a$? By Lemma~\ref{lem:abelian} the source cannot be any
commutative game-value monoid.
\end{question}

\section{No-go theorems for the middle tier}\label{sec:nogos}

This section records the no-go results of the structured attack
(2026-06-10) on the middle tier. Architectures covered: commutative
provenance, affine and orthogonal symmetry, bounded $B$-oracle access,
$B$-local flip rules, and $E$-translation dynamics. Verification scopes are
stated inline; the supporting probes live in \texttt{experiments/gold/}.
(Theorem~\ref{thm:A}, coin-turning linearity, opened the series in
Section~\ref{sec:linear}.)

\begin{nogo}[commutative squaring collapse; standard math]\label{thm:B}
In any commutative monoid, squaring is an endomorphism, so any quadratic
datum read through homomorphic coordinates has polarization identically
zero. In particular the extraspecial squaring map $Q$ (for $B\neq0$) cannot
live in, map into, or pull back through any commutative game monoid. This
covers misère quotients, turn-parity counters, and all coin-turning
configuration sums.
\end{nogo}

\begin{proof}
Immediate from Lemma~\ref{lem:abelian}: commutativity gives
$(mn)^2=m^2n^2$, so the squaring form is additive and $B\equiv0$.
\end{proof}

\begin{nogo}[group misère quotients are tiny; conditional]\label{thm:C}
Assume \cite[Thm.~6.4]{PS} with its regularity hypothesis. If the misère
indistinguishability quotient $\mathcal{Q}(\mathcal{A})$ of a set of
impartial games is a group, then $|\mathcal{Q}(\mathcal{A})|\leq 2$.
\end{nogo}

\begin{proof}[Proof sketch]
Doubling lemma: in a group quotient every element is cancellable, and any
position of Grundy value $2$ doubles to a misère-$P$ position,
contradicting Grundy-distinctness from $0$; the quoted structure theorem
then pins the quotient to at most $\{1,a\}$.
\end{proof}

\begin{remark}[scope and a probe caveat]\label{rem:thmC}
Theorem~\ref{thm:C} is conditional on the Plambeck--Siegel structure
theorem as cited; the regularity hypothesis is quoted from the repository's
\path{experiments/misere_kernel.py} docstring, and the JCTA 2008 paper has
\emph{not} been independently re-verified within this program (closing that
gate is a listed next move). Two consequences for the probe map: the
$(\mathbb{Z}/2)^k$ ($k\geq2$) target of \texttt{examples/octal\_hunt.rs} is
empty a priori -- its empirical negatives (quotient orders
$2,6,10,12,14$) are a theorem, not a sampling accident -- and any future
``hit'' from that sweep would be a labeling artifact, since the helper
\texttt{p\_set\_as\_f2} checks only that the subset labeling is a bijection,
not that it is a monoid homomorphism. The hunt's target should be reframed
accordingly.
\end{remark}

For the remaining results, fix a nondegenerate quadratic form $Q$ on
$V=\F_2^{2r}$, $r\geq2$, with polar form $B$ and zero set
$\Sigma=\{Q=0\}$, and write
$\AO(Q)=\{x\mapsto gx+c: g\in\Sp(B),\ Q(gx+c)=Q(x)\ \forall x\}$ for the
affine stabilizer appearing below.

\begin{nogo}[affine symmetry budget]\label{thm:D}
$\Stab_{\AGL(V)}(\Sigma)=\AO(Q)$; it contains no nontrivial pure
translation (for a nonsingular core, $c\in\rad(B)=0$), and its pure-linear
part is exactly $\Orth(Q)$. Consequently every uniform rule whose move digraph
admits even one affine automorphism outside $\AO(Q)$ has $P$-set
$\neq\Sigma$.
\end{nogo}

\noindent\emph{Verification.} Exhaustive over all $322{,}560$ affine maps
of $\F_2^4$ at $m=4$, both Arf classes
(\path{experiments/gold/skeptic_nogo_check.py}); the linear-part
identification is standard math (Witt extension).

\begin{nogo}[the symmetry dial has no middle]\label{thm:E}
$\Orth(Q)$ is maximal in $\Sp(B)$, and the $\Orth(Q)$-orbitals on $V\times V$
are exactly the Gram classes $(Q(u),Q(w),B(u,w))$ together with the
degeneracy flags (Witt extension; see \cite[\S8]{EKM08}). Hence a rule
invariant under any group strictly larger than $\Orth(Q)$ is dead by
Theorem~\ref{thm:D} (its symmetry reaches outside $\AO(Q)$), while a rule
fully $\Orth(Q)$- or $\Aut(E)$-equivariant factors extensionally through
pointwise $Q$- and $B$-evaluation. There is no invariance class strictly
between dead (too coarse) and Gram-factoring (an evaluator): $E$-naturality
must be \emph{covariance}, not invariance.
\end{nogo}

\noindent\emph{Verification.} Maximality verified at $m=4$: every one of
the $648$ (resp.\ $600$) elements of $\Sp(4,2)$ outside $\Orth^{+}(4,2)$
(resp.\ $\Orth^{-}(4,2)$) generates $\Sp(4,2)$ together with the orthogonal
group; orbitals enumerated exhaustively at $m=4$
(\texttt{experiments/gold/ao\_orbitals.py}).

\begin{nogo}[$B$-oracle lower bound]\label{thm:F}
Access model: legality of a move is decided from XOR combinations of the
two positions and $t$ fixed constant vectors, with oracle access to $B$
only. If $t\leq 2r-3$, then for \emph{every} quadratic refinement $Q'$ of
$B$ simultaneously there is a nonzero $Q'$-singular vector in the
complement of the constant span; the transvection it induces is a digraph
automorphism lying outside $\Orth(Q')$, and Theorem~\ref{thm:D} kills the
$P$-set. The bound is tight: at $t=2r-2$ with anisotropic complement the
argument fails. Hence the Gold diagonal framing $q_i=Q_a(e_i)$, which
supplies $2r$ frame constants, is within two dimensions of
information-theoretically forced.
\end{nogo}

\noindent\emph{Verification.} Transvection step and the exact escape
boundary $t=2r-2$ machine-checked at $m=4$; end-to-end spot check on $50$
random $t=1$ $B$-oracle rules, retrograde-solved
(\path{experiments/gold/skeptic_nogo_check.py},
\path{nogo_verify.py}).

\begin{nogo}[$B$-local flip rules]\label{thm:G}
Any rule of the form ``flip $d$ at $v$ iff $f(d,B(v,d))$'' is automatically
undirected, because $B$ is alternating:
$f(d,B(v+d,d))=f(d,B(v,d)+B(d,d))=f(d,B(v,d))$. For undirected loopy games
the outcome partition collapses: $\mathrm{Win}=\varnothing$,
$\mathrm{Loss}=$ the isolated vertices $=$ an affine flat, and
$\mathrm{Loss}\cup\mathrm{Draw}\neq\Sigma$ for $r\geq2$ (the count
$|\Sigma|=2^{r-1}(2^{r}+\varepsilon)$ has an odd factor $>1$). This covers
the Loss, Draw, and $\mathrm{Loss}\cup\mathrm{Draw}$ targets at once. The
symmetric-$B$ rule of \texttt{examples/loopy\_quadric.rs} is the special
case $f(d,\beta)=\beta$; its $(4,1)$ ``hit'' is the radical coincidence of
Section~\ref{sec:benches}.
\end{nogo}

\noindent\emph{Verification.} Undirectedness is the two-line identity
above (standard math); the loopy outcome collapse including the Draw-set
gap is machine-checked at $m=4$
(\path{experiments/gold/skeptic_nogo_check.py}).

\begin{nogo}[$E$-translation no-go]\label{thm:H}
For $r\geq2$, no left-$E$-equivariant rule on $E=E_Q$ with moves
$g\to gn$, $n\in N$ for a fixed nonempty $N\subseteq E$, has $P$-set equal
to the involution locus $I=\{g:g^2=1\}=\pi^{-1}(\Sigma)$.
\end{nogo}

\begin{proof}
$I\cdot I=E$: for every $v\neq0$, the count
$\#\{x:Q(x)=0=Q(x+v)\}=2^{2r-2}+(-1)^{\Arf Q}2^{r-1}$ is positive for
$r\geq2$ (character sum), so every $n\in N$ factors as $n=gh$ with
$g,h\in I$. Then $g\in I$, and $g\cdot n=g(gh)=g^2h=h\in I$: the move
$g\to gn$ is an edge from $I$ to $I$. A normal-play $P$-set admits no
$P$-to-$P$ edge, so $P\neq I$.
\end{proof}

\noindent\emph{Verification.} Cross-checked on the bent
$(8,1,\lambda=2)$ extraspecial group
(\path{experiments/gold/extraspecial_core.py}).

\subsection{Normal-play rigidity and the synthesis}

The strongest constraint is not any single architecture kill but a rigidity
statement inside a declared access class. Call a rule \emph{in-class} if the
quadratic refinement is quarantined behind a weight-bounded oracle (the
$(w_0,c)$ metering of Definition~\ref{def:criterion} below) and the rule is
\emph{refinement-uniform}: it realizes the target for every refinement $q$
of $B$ in the declared family, not just the Gold one.

\begin{theorem}[normal-play rigidity, in-class]\label{thm:rigidity}
Under the bounded-framing oracle model with refinement uniformity, every
legal edge between positions of Hamming weight $>c\,w_0$ must flip $Q$ --
in normal, misère, and loopy-Loss semantics alike. Consequently every
in-class realizer of $\Sigma$ in those semantics is a \emph{$Q$-alternating
ender}: play telescopes $Q$, and the outcome is the parity of $Q(x)$.
\end{theorem}

\begin{proof}[Proof sketch]
The case $Q(v)=Q(w)=0$ of a non-flipping edge is forbidden directly by the
outcome structure in all three readings. The case $Q(v)=Q(w)=1$ is reduced
to it by an adversary substitution: replace $q$ by $q'=q+\ell$ with $\ell$
a linear form annihilating every queried functional but not the dense
indicators; refinement uniformity transports the rule to $q'$, where the
edge becomes a forbidden $0$-to-$0$ edge.
\end{proof}

\begin{theorem}[no live middle, relative to the quarantine]\label{thm:nolivemiddle}
Let $Q$ have nonsingular core of rank $2r\geq4$. No fixed uniform
single-board rule satisfying
\begin{enumerate}[label=(\Alph*),leftmargin=*,itemsep=0pt]
\item[(S)] no affine digraph automorphism outside $\AO(Q)$,
\item[(A)] legality from $B$-oracle queries at $t\leq 2r-3$ frame constants,
\item[(C)] no commutative squaring route for the provenance of $Q$,
\item[(E)] no left-$E$-equivariant structure on the extraspecial group,
\item[(R)] bounded-framing oracle access with refinement uniformity,
\end{enumerate}
has $P$/Loss/Draw target equal to $\Sigma$ with any bulk strategic
content: under (R), every legal edge between dense positions flips $Q$, all
strategic content is confined to a Hamming ball of radius $c\,w_0$, and no
bulk-decision-live realizer exists.
\end{theorem}

\begin{remark}[honest scope: five escape hatches]\label{rem:hatches}
Theorem~\ref{thm:nolivemiddle} is relative to its hypotheses, and each
hypothesis is an open experimental window:
\begin{enumerate}[leftmargin=*,itemsep=0pt]
\item \emph{Loopy-Draw semantics}: the substitution contradiction of
Theorem~\ref{thm:rigidity} does not fire (Draw-to-Draw edges are legal).
\item \emph{$t\geq 2r-2$ with anisotropic complement}: escapes
Theorem~\ref{thm:F} at the symmetry level.
\item \emph{Frobenius-aware access}: the Galois group lies inside
$\Orth(Q_a)$, so all symmetry methods are silent;
Theorem~\ref{thm:F}'s model explicitly excludes Frobenius-twisted queries.
\item \emph{Non-quarantined rules using the game-native $\wp(w)$ diagonal
source} (Proposition~\ref{prop:wp}): once $q$ is game-built, refinement
uniformity holds trivially and class (R) says nothing.
\item \emph{Rank $r=1$ and radical-anisotropic degenerate layers}: the
character-sum count in Theorem~\ref{thm:H} vanishes, and small extraspecial
groups can be realized.
\end{enumerate}
The quarantine hypothesis (R) is a choice with a justification: without
quarantining the form-identifying datum, the descent evaluator is
decision-live everywhere and no purely complexity-based no-go is possible.
Whether game-built diagonal sources such as $\wp(w)$ count as ``probing''
or ``feeding in'' the form is the central open definitional question of
this program.
\end{remark}

\section{A formalized naturality criterion}\label{sec:criterion}

Two natural first attempts at ``naturality'' fail for structural reasons.
\emph{Equivariance fails}: by Theorem~\ref{thm:E} there is no invariance
class strictly between symmetry groups that kill the $P$-set and symmetry
groups that force extensional factoring through $Q$- and $B$-evaluation.
\emph{Vocabulary quarantine fails}: the Gold diagonal
$q_i^{(m,1)}=\Tr(e_i^3)$ is expressible as $\Tr(\wp(w)e_i)$ using only XOR,
nim-product, and trace (Proposition~\ref{prop:wp}) -- the licensed operation
chain itself -- so no principled line separates ``computing $q$'' from
``computing $Q$'' on vocabulary alone. The criterion below replaces both
with access metering plus a strategic-content axiom.

\begin{definition}[uniform local rules; the N-axioms]\label{def:criterion}
Fix the coin frame $\{e_i\}$ (the Grundy basis $g(n)=2^n$, game-natural).
\emph{Public data}: the frame, $\oplus$, nim-product, Frobenius, and the
full alternating form $B$ (game-built via Turning Corners, Frobenius, and
trace). \emph{Quarantined data}: the quadratic refinement, presented as a
weight-bounded oracle $O_q(z)=Q_q(z)$ for $\wt(z)\leq w_0$.
A \emph{uniform local rule} satisfies:
\begin{enumerate}[leftmargin=*,itemsep=0pt]
\item[\textbf{F1}] (\emph{$q$-blind statics}) position set, loading map,
and terminal set are computed without oracle access;
\item[\textbf{F2}] (\emph{metered access}) each legality bit is decided
with at most $c$ adaptive oracle queries, each at a point of Hamming weight
at most $w_0$, with $(w_0,c)$ constants independent of $m$;
\item[\textbf{F3}] (\emph{declared semantics}) one canonical outcome
labeling (normal, misère, loopy, or $q$-blind terminal payoff).
\end{enumerate}
It is a \emph{natural realizer} if additionally:
\begin{enumerate}[leftmargin=*,itemsep=0pt]
\item[\textbf{N1}] (\emph{torsor uniformity -- anti-evaluator}) for every
$m$, every $B$ in the declared family, and every refinement $q$ of $B$:
$\mathrm{outcome}(\iota(x))=\text{target}$ iff $Q_q(x)=0$;
\item[\textbf{N2}] (\emph{locality budget}) the $(w_0,c)$ metering of F2;
\item[\textbf{N3}] (\emph{strategic relevance -- anti-clock}) for every
instance $(m,B,q)$ there is a reachable position that is simultaneously
\emph{outcome-critical} (it has legal moves to different outcome classes)
and \emph{form-live} (its outcome differs under some refinement $q'$).
\end{enumerate}
A rule satisfying the outcome-critical clause of N3 at some reachable
position is called \emph{decision-nondegenerate}.
\end{definition}

\begin{theorem}[Tier-3 exclusion under N1+N2]\label{thm:tier3excl}
For any position $x$ with $\wt(x)>c\,w_0$, the at most $c$ queried
functionals span only weight-$\leq c\,w_0$ directions, so there exist
refinements $q,q'$ agreeing with every oracle answer but differing on
$Q(x)$. Hence no in-class rule's move predicate factors through $Q(x)$ at
dense positions -- by theorem, not stipulation.
\end{theorem}

\begin{theorem}[clock completeness; machine-verified]\label{thm:clock}
N1+N2 alone, without N3, admit pure transport clocks in every semantics,
including plain normal play: the \emph{pending-marker clock} -- positions
$(x,\varepsilon,i)$ with invariant $J=Q(x)\oplus\varepsilon$, $(1,1)$-local,
zero outcome-critical positions -- satisfies N1+N2 exactly for every
characteristic-$2$ form. Hence N3 is load-bearing: it carries the entire
residual content of ``natural''.
\end{theorem}

\begin{remark}[N3 is stated to be attacked]\label{rem:N3}
The \emph{escape-edge} construction -- a clock plus one dominated $q$-blind
gadget edge from every nonzero position -- passes N1, N2, and N3 (and a
density-strengthened N3$^{+}$) while being morally a clock
(\path{experiments/gold/skeptic_escape_edge_attack.py}). The natural
repair, a dominance-pruned or two-game version of N3, runs into the fact
that two-game criticality (critical in both the $q$- and $q'$-games with
differing outcomes) is unsatisfiable in two-class outcome semantics. The
anti-clock axiom therefore remains an open definitional problem; N3 is the
best current formulation, recorded so that it can be attacked further.
\end{remark}

\paragraph{Calibration (implemented and tested at the stated instances).}
Under F1--F3 + N1--N2: the T2 spin-flip rule of
Section~\ref{sec:constructions} fails N3 (zero outcome-critical positions,
verified at $(8,1)$ and $(8,2)$); the pending-marker clock fails N3;
\textsc{echo}-ko passes N3 (at least $160$ mistake-states at each verified
exact $m=4$ instance); an order-$4$-alphabet loopy candidate from the run
fails N2 (its alphabet is a $2^m$-point $Q$-evaluation); an
unwinding-game candidate from the run fails N1 (isotropic frames break
exactness); and the repository's Tier-3 circuit and
\texttt{interactive\_kernel} Rule~2 fail N2 (dense-point form access).
The candidates named here live in \texttt{experiments/gold/}.

\section{Constructions and near-misses}\label{sec:constructions}

\subsection{T2: the width-2 spin-flip rule (exact, but a clock)}

\begin{proposition}[T2 realizes every char-2 quadric; implemented and
tested]\label{prop:t2}
On positions $x\in\F_2^n$, let the rule be: turn over a set $d$ of
coins with $\wt(d)\in\{1,2\}$, leading coin heads-to-tails, legal iff
$B(x,d)\oplus Q(d)=1$. Then for every characteristic-$2$ quadratic form
$Q$, the normal-play $P$-set is exactly $\{Q=0\}$.
\end{proposition}

\begin{proof}[Proof sketch (blocking lemma)]
Every legal move flips $Q$ (by the polarization identity, since
$Q(x+d)=Q(x)+Q(d)+B(x,d)$ and legality forces $Q(d)+B(x,d)=1$), so it
suffices that every position with $Q(v)=1$ has a legal descent. If all
descents from $v$ fail, then $B(v,e_i)=q_i$ for all $i\in\mathrm{supp}(v)$
and $B\equiv0$ on $\mathrm{supp}(v)\times\mathrm{supp}(v)$; expanding
$Q(v)$ in the basis then forces $Q(v)=0$, a contradiction.
\end{proof}

\noindent\emph{Verification.} Checked at $(8,1)$, $(16,1)$, the bent
$(8,1,\lambda=2)$, and $20$ random refinements of each $B$ (the
``refinement uniformity certificate'';
\texttt{experiments/gold/witness\_test.py} and companions). For $a=1$ the
rule constants $q_i$ are game-native via Proposition~\ref{prop:wp}; for
$a\geq2$ they are fed in, and game-nativity is open.

T2 is, however, a forced clock: every legal move flips $Q$, play telescopes
$Q(x)$ as a parity count, every position's options share one outcome class,
and no position has a losing move. By Theorem~\ref{thm:rigidity} this is
unavoidable for \emph{all} in-class normal-play realizers, not a defect of
T2 specifically. T2 and the pending-marker clock pin the boundary: within
every tested bounded-access quarantine, exact realizers exist but are
forced clocks.

\subsection{\textsc{echo}-ko: charge counting on the extraspecial cocycle}

The one known exact realizer that is \emph{not} a clock works one level up,
on the extraspecial cocycle itself.

\begin{quote}
\textbf{Rule (\textsc{echo}-ko).} Positions are indexed by
$x\in\F_{2^m}$; the coins of $x$ must each be touched twice; players
alternate single touches; each touch of $e_i$ updates a charge
$\sigma\mathrel{\oplus{=}}c(\text{open-set},e_i)$, where $c$ is the
triangular cocycle of the extraspecial extension
($c(v,v)=Q_a(v)$, $c(u,v)\oplus c(v,u)=B(u,v)$,
cf.\ Lemma~\ref{lem:extdict}(ii)); a one-move ko forbids immediately
re-touching; the player completing the board wins iff the terminal charge
is $\sigma=1$.
\end{quote}

The rule is center-reading (not center-blind), noncommutative via move
order (not turn count), and -- at every verified exact instance --
decision-nondegenerate. For $a=1$ its constants are game-built via
Proposition~\ref{prop:wp}, with no per-instance $Q$-evaluation.

\paragraph{Provenance and results (implemented and tested).}
A first-round solver had a bug (the memoization key omitted the accumulated
charge parity); the corrected solver was validated against explicit tree
enumeration. Corrected results
(\texttt{experiments/gold/echo\_charge\_probe.py} and companions):
\begin{center}
\begin{tabular}{llcc}
\toprule
instance & core & agreement & decision-live \\
\midrule
$m=4$, bent $\lambda\in\{2,12\}$ & rank $4$ & $16/16$ exact & yes ($\geq160$ mistake-states) \\
$(8,2)$, $\lambda=1$ & rank $4$ & $255/256$ (miss $x=224$) & yes \\
$(8,1)$, $\lambda=1$ & rank $6$ & $228/256$ & yes \\
$(8,1)$, bent $\lambda=2$ & rank $8$ & $212/256$ & yes \\
\bottomrule
\end{tabular}
\end{center}
Agreement grades by polar rank. A bounded-memory blocker conjecture (for
fixed ko-window $w$, adversarial unlinking wins on supports $k\geq w+2$ at
rank $\geq6$) was formulated on the \emph{invalid} round-1 data and has not
been re-examined against the corrected solver; it is recorded as a
hypothesis to re-test, not a finding.

\subsection{\textsc{echo}-\textsc{fifo}+dummy: an unverified exactness claim}

A second-round variant replaces the one-move ko with a \textsc{fifo}
close-discipline (only the longest-open coin may be closed) and adjoins one
neutral tempo coin. The run claims full $m=8$ exactness
($391{,}680$ checks across all $765$ scaled Gold forms, including rank $6$
and rank $8$) via a decomposition-plus-isomorphism-caching solver that was
validated only at $m=4$, not directly at $m=8$.

\textbf{This claim is unverified pending adversarial review and is not
treated as established anywhere in this note.} A fresh direct stateful
solver -- no decomposition extrapolation, full state in the memo key
including $\sigma$, validated against explicit tree enumeration -- has not
been run at $m=8$. It is the single most load-bearing unverified result of
the program, and re-deriving it (or refuting it) is the top-ranked next
move.

\subsection{The decisive experiment}

The pre-registered experiment for the primary candidate: a corrected-solver
sweep over the finite \textsc{echo} rule family at $m=8$, validated by a
direct stateful solver.
\begin{itemize}[leftmargin=*,itemsep=0pt]
\item \emph{Benchmarks}: $(8,1,\lambda=1)$ rank $6$; bent
$(8,1,\lambda=2)$ rank $8$; $(8,2,\lambda=1)$ rank $4$; the full $m=4$
family as regression.
\item \emph{Axes}: ko-memory window $w\in\{1,2,3\}$; pass semantics
(clears-ko / forbidden / loses); single-coin plus pair touches (the
tartan-companion axis); both orientations.
\item \emph{Success criterion}: zero misses across all tested
$(m,a,\lambda)$ triples, decision-nondegenerate throughout, and surviving a
torsor sweep ($\geq20$ stratified refinements of each $B$, confirming
refinement uniformity rather than single-refinement luck).
\end{itemize}
A \textsc{confirm} outcome would be the first genuine Tier-2 witness
(pending recasting the charge readout into normal/misère/loopy semantics
and the even-$a$ diagonal lemma). A \textsc{kill} outcome -- rank-graded
decay persisting across every family member -- would put the
bounded-memory blocker conjecture on valid data and close the line for
bounded-memory architectures. Either outcome is informative, and the sweep
is feasible in minutes per candidate on the existing infrastructure.

\section{Status and next moves}\label{sec:status}

\paragraph{Status of the trichotomy.}
Tier 1 ($\Sp(B)$-blind) is dead by Proposition~\ref{prop:nogo} and
Theorems~\ref{thm:D}--\ref{thm:G}. Tier 3 (evaluation) is live and even
decision-nondegenerate, but tautological. Tier 2 under the bounded-framing
quarantine is inhabited \emph{only by clocks}
(Theorems~\ref{thm:rigidity}, \ref{thm:clock}): N1+N2 are satisfiable, and
every known in-class normal-play realizer violates N3. The live question is
whether a decision-nondegenerate Tier-2 rule -- N1, N2, and N3
simultaneously -- exists. \textsc{echo}-ko is the leading candidate; the
\textsc{echo}-\textsc{fifo}+dummy exactness claim is the most important
unverified result in the program.

\paragraph{Ranked next moves.}
\begin{enumerate}[leftmargin=*]
\item \emph{Adversarial review of \textsc{echo}-\textsc{fifo}+dummy}
(highest leverage): build a fresh direct stateful solver and run it on the
$m=8$ benchmarks plus $\geq20$ stratified $\lambda$ from the claimed sweep;
cross-validate the decomposition lemma directly at $m=8$; measure
N3-conjunction and bulk-liveness.
\item \emph{Pre-registered \textsc{echo}-ko family sweep}
(Section~\ref{sec:constructions}): \textsc{confirm} resolves the problem
constructively pending minor lemmas; \textsc{kill} grounds the
bounded-memory blocker on valid data.
\item \emph{Mechanism reconciliation}: derive the \textsc{fifo}+dummy
reduced-game facts and re-run the extraspecial support-counting argument
under \textsc{fifo} discipline; feeds both items above.
\item \emph{Frobenius-aware access enumeration}: finite enumeration of
Frobenius-equivariant legality predicates at $m=4,8$ -- the one access
window (escape hatch 3 of Remark~\ref{rem:hatches}) where both the symmetry
and oracle methods are provably silent.
\item \emph{Cheap gates}: (a) verify the regularity hypothesis of
\cite[Thm.~6.4]{PS} against the published paper (load-bearing for
Theorem~\ref{thm:C}); (b) enumerate conjugation-move rules on $E$ (the
involution locus is conjugacy-invariant, so Theorem~\ref{thm:H}'s
left-translation kill does not apply); (c) prove or refute the even-$a$
$\wp$-preimage family (Section~\ref{sec:framing}); (d) exhaust the board-8
case of the \textsc{fifo} parity-pinning conjecture.
\end{enumerate}

\section{Validation}\label{sec:validation}

The claims above are backed by implementation checks, not by a finished
proof of the open game statement.

\begin{itemize}[leftmargin=*]
\item \texttt{cargo test} exercises the scalar arithmetic, Clifford product,
Arf computation, quadratic fitting, misere machinery, and related invariant
modules.
\item \path{experiments/trace_form_arf.py} checks the Gold rank formula
(Table~\ref{tab:gold}).
\item \path{experiments/gold_form_from_games.py} and
\path{experiments/tartan_bilinear.py} rebuild the Gold form and polar form
from literal Turning-Corners products in small fields.
\item \path{experiments/arf_win_bias.py} checks the zero-count formula
(Table~\ref{tab:bias}).
\item The \texttt{misere\_quotient}, \texttt{octal\_hunt},
\texttt{interactive\_kernel}, \texttt{loopy\_quadric}, and
\texttt{bent\_route} examples implement the route probes of
Section~\ref{sec:benches}.
\item \path{experiments/gold_family_survey.py} checks the scaled-component
bent counts; \path{experiments/framing_obstruction.py} the diagonal-framing
pattern; \path{experiments/misere_kernel.py} the $R_8$ kernel reading.
\item \texttt{experiments/gold/} carries the probes of the 2026-06-10
parallel run backing Sections~\ref{sec:nogos}--\ref{sec:constructions}:
exhaustive symmetry checks at $m=4$, the $B$-oracle and flip-rule kills,
the extraspecial and \textsc{echo} solvers, and the criterion calibration.
These are machine-generated research probes, run under the project
virtualenv, and are \emph{not} CI-maintained; each result cited from them
states its verification scope inline above.
\end{itemize}

Two probe-hygiene caveats from the run remain open in the repository and
are recorded in Remark~\ref{rem:thmC}: the \texttt{octal\_hunt} target
should be reframed (its group-quotient target is empty under
Theorem~\ref{thm:C}), and \texttt{p\_set\_as\_f2} should verify that its
subset labeling is a homomorphism.

\section{Conclusion}

The meaningful content is layered. At the base: game-operation-built
characteristic-$2$ trace forms inside a nimber Clifford backend, with a
precise conditional interpretation of Arf as a $P$-set bias. Above it: the
linear ceiling (Theorem~\ref{thm:A}, with the lexicodes as its rich solved
floor), the extraspecial identification of the missing quadratic datum with
a noncommutative central extension (Lemmas~\ref{lem:extdict}--\ref{lem:abelian}),
and a no-go program (Theorems~\ref{thm:B}--\ref{thm:H},
\ref{thm:rigidity}--\ref{thm:nolivemiddle}) that empties every symmetric,
commutative, and oracle-bounded architecture of bulk strategic content. At
the boundary: exact realizers exist in every tested quarantine class but
are forced clocks; one construction (\textsc{echo}-ko) is exact on the
$m=4$ bent instances, degrades gracefully with polar rank, and stays
decision-live; and a formalized criterion (N1--N3) reduces
``natural'' to one load-bearing, still-attackable axiom. The draft does not
solve the game-semantics problem. It narrows it to a named question:
\emph{does any fixed uniform rule satisfy N1, N2, and N3 simultaneously on
a Gold quadric of core rank $\geq6$?}

\section*{Acknowledgements}
The \texttt{ogdoad} library, the experiments, and this draft were developed
with LLM assistance; Sections~\ref{sec:nogos}--\ref{sec:status} synthesize
the output of a structured parallel research run (2026-06-10) whose probes
are committed under \texttt{experiments/gold/}. Results from that run are
reported with their verification scope, and one construction claim is
explicitly flagged unverified. Mathematical direction, framing, and
responsibility for claims remain the author's.

\begin{thebibliography}{9}
\bibitem{Arf41} C.~Arf, \emph{Untersuchungen \"uber quadratische Formen in
K\"orpern der Charakteristik~2, I}, J.\ Reine Angew.\ Math.\ \textbf{183}
(1941), 148--167.
\bibitem{Bouton} C.~L.~Bouton, \emph{Nim, a game with a complete mathematical
theory}, Ann.\ of Math.\ \textbf{3} (1901/02), 35--39.
\bibitem{WW} E.~R.~Berlekamp, J.~H.~Conway, R.~K.~Guy, \emph{Winning Ways for
Your Mathematical Plays}, 2nd ed., A~K Peters, 2001--2004.
\bibitem{Conway76} J.~H.~Conway, \emph{On Numbers and Games}, Academic Press,
1976.
\bibitem{ConwaySloane86} J.~H.~Conway, N.~J.~A.~Sloane, \emph{Lexicographic
codes: error-correcting codes from game theory}, IEEE Trans.\ Inform.\
Theory \textbf{32} (1986), 337--348.
\bibitem{Dickson01} L.~E.~Dickson, \emph{Linear Groups, with an Exposition of
the Galois Field Theory}, Teubner, 1901.
\bibitem{EKM08} R.~Elman, N.~Karpenko, A.~Merkurjev, \emph{The Algebraic and
Geometric Theory of Quadratic Forms}, AMS Colloquium Publications
\textbf{56}, American Mathematical Society, 2008.
\bibitem{Gold68} R.~Gold, \emph{Maximal recursive sequences with $3$-valued
recursive cross-correlation functions}, IEEE Trans.\ Inform.\ Theory
\textbf{14} (1968), 154--156.
\bibitem{Gorenstein80} D.~Gorenstein, \emph{Finite Groups}, 2nd ed.,
Chelsea Publishing, New York, 1980.
\bibitem{LN} R.~Lidl, H.~Niederreiter, \emph{Finite Fields}, 2nd ed.,
Cambridge University Press, 1997.
\bibitem{Ovsienko16} V.~Ovsienko, \emph{Real Clifford algebras and quadratic
forms over $\F_2$: two old problems become one}, The Mathematical Intelligencer
\textbf{38} (2016); arXiv:1601.07664.
\bibitem{PS} T.~E.~Plambeck, A.~N.~Siegel, \emph{Misere quotients for impartial
games}, J.\ Combin.\ Theory Ser.\ A \textbf{115} (2008), 593--622.
\bibitem{Quillen71} D.~Quillen, \emph{The mod~$2$ cohomology rings of
extra-special $2$-groups and the spectra of quadratic forms}, Math.\ Ann.\
\textbf{194} (1971), 197--212.
\bibitem{Taylor} D.~E.~Taylor, \emph{The Geometry of the Classical Groups},
Sigma Series in Pure Mathematics~\textbf{9}, Heldermann Verlag, 1992.
\bibitem{Winter72} D.~L.~Winter, \emph{The automorphism group of an
extraspecial $p$-group}, Rocky Mountain J.\ Math.\ \textbf{2} (1972),
159--168.
\end{thebibliography}

\end{document}
